\documentclass[12pt]{article}
\usepackage{latexblog}

% ============================================================
% Blog Metadata
% ============================================================
\blogtitle{在Android 上使用 OpenCV}
\blogdate{2017-11-13}
\blogtags{Android, 移动端, 闲话编程, OpenCV}
\bloglang{zh}
\blogsummary{}

\begin{document}
\makeblogtitle

如题，本文将记录如何在安卓上调用OpenCV。

\section{导入OpenCV动态库}\label{ux5bfcux5165opencvux52a8ux6001ux5e93}

首先当然是下载OpenCV for Android了，然后使用Android Studio创建一个工程并勾选C++ support。

然后，把OpenCV-android-sdk里面的native目录拷贝到工程中，例如app/opencv这个目录，需要修改以下文件：

\begin{itemize}
\tightlist
\item
  app/build.gradle
\end{itemize}

\begin{Shaded}
\begin{Highlighting}[]
\NormalTok{android }\OperatorTok{\{}
\OperatorTok{....}
\FunctionTok{sourceSets} \OperatorTok{\{}
\NormalTok{        main }\OperatorTok{\{}
\NormalTok{            jniLibs}\OperatorTok{.}\NormalTok{srcDirs }\OperatorTok{=} \OperatorTok{[}\StringTok{\textquotesingle{}opencv/libs\textquotesingle{}}\OperatorTok{]}
        \OperatorTok{\}}
    \OperatorTok{\}}
\OperatorTok{\}}
\end{Highlighting}
\end{Shaded}

这样做的目的是为了打包的时候能自动将opencv/libs/\{arch\}/libopencv\_java3.so这个文件打包到我们的apk中。

\begin{itemize}
\tightlist
\item
  app/CMakeList.txt
\end{itemize}

\begin{verbatim}
set(opencv "${CMAKE_SOURCE_DIR}/opencv")
include_directories(${opencv}/jni/include)
add_library(libopencv_java3 SHARED IMPORTED )
set_target_properties(libopencv_java3 PROPERTIES
                      IMPORTED_LOCATION "${opencv}/libs/${ANDROID_ABI}/libopencv_java3.so")

target_link_libraries( # Specifies the target library.
                       native-lib
                       libopencv_java3
                       ${log-lib} )
\end{verbatim}

这里把opencv作为动态库链接到工程中，并添加了包含目录，否则在编译cpp的时候会找不到opencv的头文件。

\section{导入OpenCV Jaba Module}\label{ux5bfcux5165opencv-jaba-module}

把opencv sdk下面的java目录作为一个module导入到工程中，并设置app依赖此module，这样就可以在工程中使用opencv提供的java接口了。我们来做一个相机：

\begin{Shaded}
\begin{Highlighting}[]
\KeywordTok{public} \KeywordTok{class}\NormalTok{ MainActivity }\KeywordTok{extends}\NormalTok{ AppCompatActivity }\KeywordTok{implements}\NormalTok{ CameraBridgeViewBase}\OperatorTok{.}\FunctionTok{CvCameraViewListener2}
\end{Highlighting}
\end{Shaded}

这里首先实现CvCameraViewListener2接口，来实现相机的处理。

\begin{Shaded}
\begin{Highlighting}[]
    \KeywordTok{private}\NormalTok{ CameraBridgeViewBase cameraView}\OperatorTok{;}

    \KeywordTok{private}\NormalTok{ BaseLoaderCallback loaderCallback }\OperatorTok{=} \KeywordTok{new} \FunctionTok{BaseLoaderCallback}\OperatorTok{(}\KeywordTok{this}\OperatorTok{)} \OperatorTok{\{}
        \AttributeTok{@Override}
        \KeywordTok{public} \DataTypeTok{void} \FunctionTok{onManagerConnected}\OperatorTok{(}\DataTypeTok{int}\NormalTok{ status}\OperatorTok{)} \OperatorTok{\{}
            \ControlFlowTok{switch} \OperatorTok{(}\NormalTok{status}\OperatorTok{)} \OperatorTok{\{}
                \ControlFlowTok{case}\NormalTok{ LoaderCallbackInterface}\OperatorTok{.}\FunctionTok{SUCCESS}\OperatorTok{:}
\NormalTok{                    cameraView}\OperatorTok{.}\FunctionTok{enableView}\OperatorTok{();}
                    \ControlFlowTok{break}\OperatorTok{;}
                \KeywordTok{default}\OperatorTok{:}
                    \KeywordTok{super}\OperatorTok{.}\FunctionTok{onManagerConnected}\OperatorTok{(}\NormalTok{status}\OperatorTok{);}
                    \ControlFlowTok{break}\OperatorTok{;}
            \OperatorTok{\}}
        \OperatorTok{\}}
    \OperatorTok{\};}

    \AttributeTok{@Override}
    \KeywordTok{protected} \DataTypeTok{void} \FunctionTok{onCreate}\OperatorTok{(}\NormalTok{Bundle savedInstanceState}\OperatorTok{)} \OperatorTok{\{}
        \KeywordTok{super}\OperatorTok{.}\FunctionTok{onCreate}\OperatorTok{(}\NormalTok{savedInstanceState}\OperatorTok{);}
        \KeywordTok{this}\OperatorTok{.}\FunctionTok{requestPermissions}\OperatorTok{();}

        \KeywordTok{this}\OperatorTok{.}\FunctionTok{getWindow}\OperatorTok{().}\FunctionTok{addFlags}\OperatorTok{(}\NormalTok{WindowManager}\OperatorTok{.}\FunctionTok{LayoutParams}\OperatorTok{.}\FunctionTok{FLAG\_KEEP\_SCREEN\_ON}\OperatorTok{);}
        \FunctionTok{setContentView}\OperatorTok{(}\NormalTok{R}\OperatorTok{.}\FunctionTok{layout}\OperatorTok{.}\FunctionTok{activity\_main}\OperatorTok{);}
        \KeywordTok{this}\OperatorTok{.}\FunctionTok{cameraView} \OperatorTok{=} \OperatorTok{(}\NormalTok{CameraBridgeViewBase}\OperatorTok{)} \KeywordTok{this}\OperatorTok{.}\FunctionTok{findViewById}\OperatorTok{(}\NormalTok{R}\OperatorTok{.}\FunctionTok{id}\OperatorTok{.}\FunctionTok{cameraView}\OperatorTok{);}
        \KeywordTok{this}\OperatorTok{.}\FunctionTok{cameraView}\OperatorTok{.}\FunctionTok{setVisibility}\OperatorTok{(}\NormalTok{SurfaceView}\OperatorTok{.}\FunctionTok{VISIBLE}\OperatorTok{);}
        \KeywordTok{this}\OperatorTok{.}\FunctionTok{cameraView}\OperatorTok{.}\FunctionTok{setCvCameraViewListener}\OperatorTok{(}\KeywordTok{this}\OperatorTok{);}
    \OperatorTok{\}}
\end{Highlighting}
\end{Shaded}

在create的时候，我们申请权限，然后设置相机view的监听为自身。

\begin{Shaded}
\begin{Highlighting}[]
    \AttributeTok{@Override}
    \KeywordTok{public} \DataTypeTok{void} \FunctionTok{onCameraViewStarted}\OperatorTok{(}\DataTypeTok{int}\NormalTok{ width}\OperatorTok{,} \DataTypeTok{int}\NormalTok{ height}\OperatorTok{)} \OperatorTok{\{}
    \OperatorTok{\}}

    \AttributeTok{@Override}
    \KeywordTok{public} \DataTypeTok{void} \FunctionTok{onCameraViewStopped}\OperatorTok{()} \OperatorTok{\{}
    \OperatorTok{\}}
\end{Highlighting}
\end{Shaded}

相机启动停止我们不需要做别的操作。

\begin{Shaded}
\begin{Highlighting}[]
    \AttributeTok{@Override}
    \KeywordTok{protected} \DataTypeTok{void} \FunctionTok{onResume}\OperatorTok{()} \OperatorTok{\{}
        \KeywordTok{super}\OperatorTok{.}\FunctionTok{onResume}\OperatorTok{();}
        \ControlFlowTok{if} \OperatorTok{(!}\NormalTok{OpenCVLoader}\OperatorTok{.}\FunctionTok{initDebug}\OperatorTok{())}
\NormalTok{            OpenCVLoader}\OperatorTok{.}\FunctionTok{initAsync}\OperatorTok{(}\NormalTok{OpenCVLoader}\OperatorTok{.}\FunctionTok{OPENCV\_VERSION\_3\_0\_0}\OperatorTok{,} \KeywordTok{this}\OperatorTok{,} \KeywordTok{this}\OperatorTok{.}\FunctionTok{loaderCallback}\OperatorTok{);}
        \ControlFlowTok{else}
            \KeywordTok{this}\OperatorTok{.}\FunctionTok{loaderCallback}\OperatorTok{.}\FunctionTok{onManagerConnected}\OperatorTok{(}\NormalTok{LoaderCallbackInterface}\OperatorTok{.}\FunctionTok{SUCCESS}\OperatorTok{);}
    \OperatorTok{\}}
\end{Highlighting}
\end{Shaded}

相机继续的时候，我们重新加载OpenCV库。

\begin{Shaded}
\begin{Highlighting}[]
    \AttributeTok{@Override}
    \KeywordTok{public} \DataTypeTok{void} \FunctionTok{onPause}\OperatorTok{()} \OperatorTok{\{}
        \KeywordTok{super}\OperatorTok{.}\FunctionTok{onPause}\OperatorTok{();}
        \ControlFlowTok{if} \OperatorTok{(}\KeywordTok{this}\OperatorTok{.}\FunctionTok{cameraView} \OperatorTok{!=} \KeywordTok{null}\OperatorTok{)}
            \KeywordTok{this}\OperatorTok{.}\FunctionTok{cameraView}\OperatorTok{.}\FunctionTok{disableView}\OperatorTok{();}
    \OperatorTok{\}}

    \AttributeTok{@Override}
    \KeywordTok{public} \DataTypeTok{void} \FunctionTok{onDestroy}\OperatorTok{()} \OperatorTok{\{}
        \KeywordTok{super}\OperatorTok{.}\FunctionTok{onDestroy}\OperatorTok{();}
        \ControlFlowTok{if} \OperatorTok{(}\KeywordTok{this}\OperatorTok{.}\FunctionTok{cameraView} \OperatorTok{!=} \KeywordTok{null}\OperatorTok{)}
            \KeywordTok{this}\OperatorTok{.}\FunctionTok{cameraView}\OperatorTok{.}\FunctionTok{disableView}\OperatorTok{();}
    \OperatorTok{\}}
\end{Highlighting}
\end{Shaded}

暂停和销毁的时候，我们把相机禁用掉。

\begin{Shaded}
\begin{Highlighting}[]
    \AttributeTok{@Override}
    \KeywordTok{public}\NormalTok{ Mat }\FunctionTok{onCameraFrame}\OperatorTok{(}\NormalTok{CameraBridgeViewBase}\OperatorTok{.}\FunctionTok{CvCameraViewFrame}\NormalTok{ inputFrame}\OperatorTok{)} \OperatorTok{\{}
\NormalTok{        Mat frame }\OperatorTok{=}\NormalTok{ inputFrame}\OperatorTok{.}\FunctionTok{rgba}\OperatorTok{();}
\NormalTok{        Core}\OperatorTok{.}\FunctionTok{rotate}\OperatorTok{(}\NormalTok{frame}\OperatorTok{,}\NormalTok{ frame}\OperatorTok{,}\NormalTok{ Core}\OperatorTok{.}\FunctionTok{ROTATE\_90\_CLOCKWISE}\OperatorTok{);}
        \ControlFlowTok{return}\NormalTok{ frame}\OperatorTok{;}
    \OperatorTok{\}}
\end{Highlighting}
\end{Shaded}

这是关键的一步，处理相机的一帧。我们队图像进行了旋转，否则图像的坐标和我们的预期的是不一致的。注意在OpenCV3.2的时候，引入了便捷的rotate函数，如果用之前的方法，可能需要flip和reverse来实现了。

\begin{Shaded}
\begin{Highlighting}[]
\KeywordTok{private} \DataTypeTok{void} \FunctionTok{requestPermissions}\OperatorTok{()} \OperatorTok{\{}
        \DataTypeTok{int}\NormalTok{ permissionCheck }\OperatorTok{=}\NormalTok{ ContextCompat}\OperatorTok{.}\FunctionTok{checkSelfPermission}\OperatorTok{(}\KeywordTok{this}\OperatorTok{,} \BuiltInTok{Manifest}\OperatorTok{.}\FunctionTok{permission}\OperatorTok{.}\FunctionTok{CAMERA}\OperatorTok{);}
        \ControlFlowTok{if} \OperatorTok{(}\NormalTok{permissionCheck }\OperatorTok{==}\NormalTok{ PackageManager}\OperatorTok{.}\FunctionTok{PERMISSION\_GRANTED}\OperatorTok{)}
            \ControlFlowTok{return}\OperatorTok{;}

\NormalTok{        ActivityCompat}\OperatorTok{.}\FunctionTok{requestPermissions}\OperatorTok{(}\KeywordTok{this}\OperatorTok{,}
                \KeywordTok{new} \BuiltInTok{String}\OperatorTok{[]\{}\BuiltInTok{Manifest}\OperatorTok{.}\FunctionTok{permission}\OperatorTok{.}\FunctionTok{CAMERA}\OperatorTok{\},}
                \DecValTok{0}\OperatorTok{);}

    \OperatorTok{\}}
\end{Highlighting}
\end{Shaded}

最后是权限的动态申请。当然了，在AndroidManifest.xml中也需要进行设置，我们直接贴代码了：

\begin{Shaded}
\begin{Highlighting}[]
\FunctionTok{\textless{}?xml} \OtherTok{version=}\StringTok{"1.0"} \OtherTok{encoding=}\StringTok{"utf{-}8"}\FunctionTok{?\textgreater{}}
\NormalTok{\textless{}}\KeywordTok{manifest} \OtherTok{xmlns:android=}\StringTok{"http://schemas.android.com/apk/res/android"}
    \OtherTok{package=}\StringTok{"com.riguz.okapia"}\NormalTok{\textgreater{}}

\NormalTok{    \textless{}}\KeywordTok{uses{-}permission} \OtherTok{android:name=}\StringTok{"android.permission.WRITE\_EXTERNAL\_STORAGE"}\NormalTok{ /\textgreater{}}
\NormalTok{    \textless{}}\KeywordTok{uses{-}permission} \OtherTok{android:name=}\StringTok{"android.permission.READ\_EXTERNAL\_STORAGE"}\NormalTok{ /\textgreater{}}
\NormalTok{    \textless{}}\KeywordTok{uses{-}permission} \OtherTok{android:name=}\StringTok{"android.permission.CAMERA"}\NormalTok{ /\textgreater{}}

\NormalTok{    \textless{}}\KeywordTok{uses{-}feature} \OtherTok{android:name=}\StringTok{"android.hardware.camera"}\NormalTok{ /\textgreater{}}
\NormalTok{    \textless{}}\KeywordTok{uses{-}feature} \OtherTok{android:name=}\StringTok{"android.hardware.camera.autofocus"}\NormalTok{ /\textgreater{}}
\NormalTok{    \textless{}}\KeywordTok{uses{-}feature} \OtherTok{android:name=}\StringTok{"android.hardware.camera.front"}\NormalTok{ /\textgreater{}}
\NormalTok{    \textless{}}\KeywordTok{uses{-}feature} \OtherTok{android:name=}\StringTok{"android.hardware.camera.front.autofocus"}\NormalTok{ /\textgreater{}}

\NormalTok{    \textless{}}\KeywordTok{application}
        \OtherTok{android:allowBackup=}\StringTok{"true"}
        \OtherTok{android:icon=}\StringTok{"@mipmap/ic\_launcher"}
        \OtherTok{android:label=}\StringTok{"@string/app\_name"}
        \OtherTok{android:roundIcon=}\StringTok{"@mipmap/ic\_launcher\_round"}
        \OtherTok{android:supportsRtl=}\StringTok{"true"}
        \OtherTok{android:theme=}\StringTok{"@style/AppTheme"}\NormalTok{\textgreater{}}
\NormalTok{        \textless{}}\KeywordTok{activity} \OtherTok{android:name=}\StringTok{".MainActivity"}\NormalTok{\textgreater{}}
\NormalTok{            \textless{}}\KeywordTok{intent{-}filter}\NormalTok{\textgreater{}}
\NormalTok{                \textless{}}\KeywordTok{action} \OtherTok{android:name=}\StringTok{"android.intent.action.MAIN"}\NormalTok{ /\textgreater{}}

\NormalTok{                \textless{}}\KeywordTok{category} \OtherTok{android:name=}\StringTok{"android.intent.category.LAUNCHER"}\NormalTok{ /\textgreater{}}
\NormalTok{            \textless{}/}\KeywordTok{intent{-}filter}\NormalTok{\textgreater{}}
\NormalTok{        \textless{}/}\KeywordTok{activity}\NormalTok{\textgreater{}}
\NormalTok{    \textless{}/}\KeywordTok{application}\NormalTok{\textgreater{}}

\NormalTok{\textless{}/}\KeywordTok{manifest}\NormalTok{\textgreater{}}
\end{Highlighting}
\end{Shaded}

参考：\href{http://blog.codeonion.com/2016/04/09/show-camera-on-android-app-using-opencv-for-android/}{Use OpenCV to show camera on android App with correct orientation}


\end{document}
